% This file is part of the modular source for Stratified Relational Workspace Theory.
% !TeX root = ../Stratified_Relational_Workspace_Theory.tex
\theorypart{XX}{Conjectures and Open Mathematical Directions}

\topic{71}{Workspace Formation Transitions}

The symbolic outline proposed a coupling threshold \(\lambda_c\). A general
theory must first state what is allowed to vary with \(\lambda\). A threshold
may concern structural coherence, causal efficacy, relational irreducibility,
or the conjunction of all certification gates. These need not occur at the
same parameter value.

\subtopic{71.1 No coordinate-free universal threshold}

\begin{proposition}[Threshold values are parameterization dependent]
\label{prop:threshold-reparam}
Suppose a model indexed by \(\lambda\) passes a declared formation certificate
exactly for \(\lambda>\lambda_c\). For every strictly increasing bijection
\(f\), the reparameterized family indexed by \(\mu=f(\lambda)\) has threshold
\[
\mu_c=f(\lambda_c).
\]
Consequently, the numerical value of a coupling threshold is not invariant
without a physically or operationally fixed parameterization.
\end{proposition}

\begin{proof}
The reparameterized realization at \(\mu\) is the original realization at
\(f^{-1}(\mu)\). It passes precisely when
\(f^{-1}(\mu)>\lambda_c\), equivalently when
\(\mu>f(\lambda_c)\).
\end{proof}

There can therefore be no model-independent numerical value of \(\lambda_c\).
The invariant object is the change in the certified organization and its local
critical behavior, not the label placed on the control parameter.

\subtopic{71.2 Gate-specific transition parameters}

For a one-parameter family \(\mathfrak R_\lambda\), define
\[
\begin{aligned}
\lambda_{\mathrm{str}}
&=
\inf\{\lambda:
 d_{\varepsilon}(\lambda)>0,
 g_{\varepsilon}(\lambda)>0\},\\
\lambda_{\mathrm{op}}
&=
\inf\{\lambda:
 \delta_{\mathrm{align}}(\lambda)
 \leq\varepsilon_{\mathrm{op}}\},\\
\lambda_{\mathrm{causal}}
&=
\inf\{\lambda:
 \eta_{\mathrm{causal}}(\lambda)
 \geq\tau_{\mathrm{causal}}\},\\
\lambda_{\mathrm{rel}}
&=
\inf\{\lambda:
 k_\varepsilon(\lambda)\geq2,
 \Delta_{\mathrm{tr}}(\lambda)
 \geq\tau_{\mathrm{tr}}\},\\
\lambda_{\mathrm{form}}
&=
\inf\{\lambda:
 \operatorname{Form}_{\Theta_{\mathrm{form}}}
 (W_\lambda)=1\}.
\end{aligned}
\]
In general,
\[
\lambda_{\mathrm{str}},
\lambda_{\mathrm{op}},
\lambda_{\mathrm{causal}},
\lambda_{\mathrm{rel}},
\lambda_{\mathrm{form}}
\]
can be distinct. A synchronized representation can appear before it becomes
causally useful, and causal coordination can appear while remaining reducible
to one centralized component.

\subtopic{71.3 A solvable nonlinear formation model \(T_{18}\) }

A minimal exact model is the scalar recurrent coherence equation
\[
\boxed{
m_{t+1}=\tanh(\lambda m_t),
\qquad
\lambda\geq0.
}
\]
The order parameter \(m_t\) represents a signed coherent mode. This model is
not a complete relational workspace. It isolates one mechanism by which a
nonzero shared mode can appear continuously from a symmetric zero state.

\begin{theorem}[Pitchfork formation transition]\label{thm:pitchfork-formation}
For the map
\[
F_\lambda(m)=\tanh(\lambda m),
\qquad
m\in[-1,1],
\]
the following hold.
\begin{enumerate}
\item If \(0\leq\lambda\leq1\), the unique fixed point is \(m=0\).
 It is globally attracting for \(\lambda<1\) and nonexpanding at
 \(\lambda=1\).
\item If \(\lambda>1\), there are exactly three fixed points
 \[
 -m_\lambda,
 \quad0,
 \quad m_\lambda,
 \qquad
 m_\lambda\in(0,1).
 \]
 The zero fixed point is unstable and the two nonzero fixed points are
 locally asymptotically stable.
\item As \(\lambda\downarrow1\),
 \[
 \boxed{
 m_\lambda
 =
 \sqrt{\frac{3(\lambda-1)}{\lambda^3}}
 +
 O\bigl((\lambda-1)^{3/2}\bigr).
 }
 \]
\end{enumerate}
\end{theorem}

\begin{proof}
A fixed point satisfies
\[
\operatorname{artanh}(m)=\lambda m.
\]
For \(m>0\), the function
\[
h(m)=\frac{\operatorname{artanh}(m)}{m}
\]
is strictly increasing from \(1\) to \(+\infty\). Hence there is no positive
solution when \(\lambda\leq1\), and exactly one positive solution when
\(\lambda>1\). Odd symmetry gives the negative solution.

The derivative is
\[
F_\lambda'(m)
=
\lambda\operatorname{sech}^2(\lambda m).
\]
At zero this equals \(\lambda\), proving stability for \(\lambda<1\) and
instability for \(\lambda>1\). At a nonzero fixed point,
\[
F_\lambda'(m_\lambda)
=
\lambda(1-m_\lambda^2)
=
\frac{(1-m_\lambda^2)\operatorname{artanh}(m_\lambda)}{m_\lambda}
<1,
\]
because
\[
\operatorname{artanh}(m)
<
\frac{m}{1-m^2}
\qquad(0<m<1).
\]
The derivative is positive, so the nonzero fixed points are locally stable.
Finally, expanding
\[
\tanh(\lambda m)
=
\lambda m-
\frac{\lambda^3m^3}{3}
+O(m^5)
\]
and solving the nonzero fixed-point equation gives the stated asymptotic.
\end{proof}

The square-root scaling is a genuine model-specific critical law. It does not
show that every workspace forms by a pitchfork bifurcation. It supplies an
exact benchmark against which more complex formation dynamics can be tested.

\subtopic{71.4 Detecting a transition from finite data}

Let \(Q(\lambda)\) be a gate-specific population order parameter and
\(\widehat Q_n(\lambda)\) its held-out estimate. If
\[
\sup_{\lambda\in\Lambda_n}
|\widehat Q_n(\lambda)-Q(\lambda)|
\leq\varepsilon_n
\]
on a finite evaluation grid \(\Lambda_n\), then every threshold crossing with
vertical margin larger than \(2\varepsilon_n\) is recovered by the same
uniform-selection argument as Theorem~\ref{thm:uniform-structure-selection}.
Near a continuous transition, the margin vanishes and sample requirements
necessarily grow. A broad, slowly varying crossover should not be reported as
a sharp phase transition without a finite-size scaling analysis.

\subtopic{71.5 General formation-transition conjecture}

For adaptive transfer systems with a fixed physical parameterization, there
may exist model classes in which
\[
\lambda<\lambda_c
\Longrightarrow
\operatorname{Form}_{\Theta_{\mathrm{form}}}(W_\lambda)=0,
\]
while
\[
\lambda>\lambda_c
\Longrightarrow
\operatorname{Form}_{\Theta_{\mathrm{form}}}(W_\lambda)=1
\]
with a stable nonzero structural mode, positive causal efficacy, and
relational irreducibility. The theorem above establishes this only for one
structural order parameter in a solvable nonlinear model. The conjunction of
all workspace gates remains model dependent.


\topic{72}{Hierarchical Stabilization Conjecture}

For appropriate recurrent systems,
\[
W^{(r+1)}
=
\Gamma_{\varepsilon_r}(W^{(r)})
\]
may converge to a stable workspace type:
\[
d_{\mathrm{id}}\!\left(
W^{(r+1)},W^{(r)}
\right)
\to0.
\]

\topic{73}{Endogenous Boundary Conjecture}

The boundary theorems in Part~XV establish recovery and stability once a score,
candidate family, and formation procedure are fixed. They do not establish
that an arbitrary adaptive system will create, represent, and maintain its own
boundary.

Let
\[
B_t^{\mathrm{sys}}
\]
be an internal membership or routing state used by the system, and let
\[
S_t^*
=
\argmax_{S\in\mathfrak S_{\mathrm{pass}}}
\mathcal B_t(S)
\]
be the population boundary defined by the external certification protocol.
The strong endogenous boundary conjecture asserts that there are adaptive
systems and nontrivial environments for which:
\[
\boxed{
\begin{aligned}
&d_B(B_t^{\mathrm{sys}},S_t^*)\to0,\\
&\operatorname{Form}_{\Theta_{\mathrm{form}}}(W_{B_t^{\mathrm{sys}}})=1,\\
&\operatorname{CE}
\left(
B_t^{\mathrm{sys}}\to
\operatorname{Routing}_{t+1}
\right)>0,\\
&\operatorname{CE}
\left(
\operatorname{BoundaryDamage}_t\to
\operatorname{RepairAction}_{t+1}
\right)>0,
\end{aligned}
}
\]
while the scaffold contains no direct encoding of \(S_t^*\).

A convincing demonstration requires four separable achievements:
\begin{enumerate}
\item \emph{discovery}: interaction history changes which boundary is learned;
\item \emph{instantiation}: an internal variable causally realizes that
 boundary;
\item \emph{use}: changing the boundary changes workspace routing and
 downstream behavior; and
\item \emph{maintenance}: detected boundary degradation causes restorative
 action.
\end{enumerate}

Theorems~\ref{thm:heldout-boundary-recovery} and
\ref{thm:boundary-drift-stability} handle external consistency of the selected
boundary. Theorems~\ref{thm:repair-selection} and
\ref{thm:self-maintenance} now handle controlled repair selection and bounded
closed-loop maintenance once a health representation, damage class, monitor,
and repair controller are given. The remaining conjectural step is endogenous
learning of the internal boundary, detector, and repair policy under open-ended
environmental change without direct scaffold encoding or external
substitution.
